% Original self-study derivation bridge, 2026-10-17.
\subsection*{Worked derivation: from unitarity to the Schr\"odinger equation}
The previous section introduced the infinitesimal time-evolution operator.
Let us fill in the intermediate algebra so that the role of
Hermiticity is visible, including all terms of first order in $\dd t$.

Take $\op U(t+\dd t,t)=\mathbf1+\op A\,\dd t+
O((\dd t)^2)$ for some operator $\op A$.
Unitarity demands
\begin{align}
\mathbf1
&=\op U^\dagger\op U\notag\\
&=\bigl[\mathbf1+\op A^\dagger\dd t+O((\dd t)^2)\bigr]
  \bigl[\mathbf1+\op A\dd t+O((\dd t)^2)\bigr]\notag\\
&=\mathbf1+(\op A^\dagger+\op A)\dd t
  +O((\dd t)^2).
\end{align}
Comparing the coefficients of $\dd t$ gives
$\op A^\dagger=-\op A$. Thus the infinitesimal
generator is anti-Hermitian. We can write
$\op A=-\ii\op H/\hbar$, where
\begin{equation}
\op H^\dagger
=(\ii\hbar\op A)^\dagger
=-\ii\hbar\op A^\dagger
=\ii\hbar\op A=\op H.
\end{equation}
For a time-dependent generator the infinitesimal relation is
\begin{equation}
\op U(t+\dd t,t)=
\mathbf1-\frac{\ii}{\hbar}\op H(t)\dd t
+O((\dd t)^2).
\end{equation}
The composition property gives
\begin{align}
\op U(t+\dd t,t_0)
&=\op U(t+\dd t,t)\op U(t,t_0)\notag\\
&=\op U(t,t_0)-\frac{\ii}{\hbar}
 \op H(t)\op U(t,t_0)\dd t+O((\dd t)^2).
\end{align}
Subtract $\op U(t,t_0)$, divide by $\dd t$ and
take $\dd t\to0$:
\begin{equation}
\boxed{\ii\hbar\partial_t\op U(t,t_0)
=\op H(t)\op U(t,t_0).}
\end{equation}
Acting on the fixed initial ket $\ket{\psi(t_0)}$
immediately yields
$\ii\hbar\partial_t\ket{\psi(t)}
=\op H(t)\ket{\psi(t)}$.

\paragraph{What changes if Hamiltonians at different times do not commute?}
At two successive short intervals $\Delta t$, the evolution is
\begin{equation}
\op U(t_2,t_0)\simeq
\left[\mathbf1-\frac{\ii}{\hbar}\op H(t_1)\Delta t\right]
\left[\mathbf1-\frac{\ii}{\hbar}\op H(t_0)\Delta t\right].
\end{equation}
Reversing their order changes the product by a term
proportional to
$[\op H(t_1),\op H(t_0)](\Delta t)^2$.
For a finite interval these ordering effects accumulate and
lead to the time-ordered exponential introduced later in this book.
The expression
$\exp[-(\ii/\hbar)\int\op H(t)\dd t]$
is valid without time ordering only if the relevant Hamiltonians
commute. In the special constant two-level example
$\op H=\tfrac{\hbar\Omega}{2}\sigma_x$,
$\sigma_x^2=\mathbf1$ gives
$\op U=\cos(\Omega t/2)\mathbf1-\ii\sin(\Omega t/2)\sigma_x$.
Direct multiplication verifies
$\op U^\dagger\op U=\mathbf1$, as expected.

\paragraph{Checkpoint.}
Repeat the first calculation using the \emph{effective} Hamiltonian
of Sec.~\ref{sec:nh-two-cavities}. Identify precisely which step
changes and use Eq.~\eqref{eq:nh-survival} to interpret the result.
